\section*{Item 5}\label{it:5}

\textbf{Enunciado}: elabore um código em .m no Matlab para projetar um conversor Boost as entradas do código são devem ser:
\begin{itemize}
	\item Tensão de entrada;
	\item Tensão de saída;
	\item Frequência de comutação;
	\item Potência na saída;
	\item Ripple de corrente no indutor e de tensão de saída em \%.
\end{itemize}

O código deve verificar se o conversor opera no MCC e MCD, caso opere no MCD deve-se retornar o valor da indutância projetada e o ciclo de trabalho. Caso opere no MCC deve-se retornar o valor da indutância projetada, capacitância de saída projetada, ciclo de trabalho no ponto de operação, corrente média no indutor, correntes máxima e mínima no indutor.

\medskip
\textbf{Resolução}: o código desenvolvido para resolução deste exercícios é mostrado no Algoritmo~\ref{lst:item5}.
	
\begin{lstlisting}[caption={Função que realiza o projeto do conversor Boost.},label={lst:item5}]
    function projeto_boost(Vg,Vo,fs,Po,rip_i,rip_v)
    %--------------------------------------------
    % Projeto de conversor Boost
    %
    % Entradas:
    %   Vg    - tensao de entrada (V)
    %   Vo    - tensao de saida (V)
    %   fs    - frequencia de chaveamento (Hz)
    %   Po    - potencia de saida (W)
    %   rip_i - ripple de corrente (% da corrente media do indutor)
    %   rip_v - ripple de tensao (% da tensao de saida)
    %--------------------------------------------
    
    Ts = 1/fs;
    
    %=========================
    % Ponto de operacao em CCM
    %=========================
    
    D       = 1 - Vg/Vo;
    Dp      = 1-D;
    R       = Vo^2/Po;
    IL      = Vg/(Dp^2*R);
    Delta_i = (rip_i/100)*IL;
    L       = Vg*D*Ts/(2*Delta_i);
    K       = 2*L/(R*Ts);
    Kcrit   = D*Dp^2;
    
    fprintf('\n');
    fprintf('-------------------------------\n');
    fprintf('Projeto do Conversor Boost\n');
    fprintf('-------------------------------\n');
    
    fprintf('D = %.4f\n',D);
    fprintf('K = %.4f\n',K);
    fprintf('Kcrit = %.4f\n',Kcrit);
    
    %================================
    % Verificacao do modo de conducao
    %================================
    
    if K >= Kcrit
    fprintf('\nModo de operacao: CCM\n\n');
    
    Delta_v = rip_v/100*Vo;
    C       = Delta_i*Ts/(8*Delta_v);
    ILmax   = IL + Delta_i/2;
    ILmin   = IL - Delta_i/2;
    
    fprintf('Indutancia      = %.4e H\n',L);
    fprintf('Capacitancia    = %.4e F\n',C);
    fprintf('Duty-cycle      = %.4f\n',D);
    fprintf('Corrente media  = %.4f A\n',IL);
    fprintf('Corrente maxima = %.4f A\n',ILmax);
    fprintf('Corrente minima = %.4f A\n',ILmin);
    else
    fprintf('\nModo de operacao: DCM\n\n');
    
    M = Vo/Vg;
    K = D^2/(M*(M-1));
    L = K*R*Ts/2;
    
    fprintf('Indutancia = %.4e H\n',L);
    fprintf('Duty-cycle = %.4f\n',D);
    
    end
    end
\end{lstlisting}

Como exemplo, utilizando valores semelhantes aos do item~\ref{it:2}, obteve-se os seguintes

\begin{lstlisting}[caption={Exemplo 1 - CCM.},label={lst:item5a}]
	>> projeto_boost(100,200,10e3,200,50,10)
	
	-------------------------------
	Projeto do Conversor Boost
	-------------------------------
	D = 0.5000
	K = 0.2500
	Kcrit = 0.1250
	
	Modo de operacao: CCM
	
	Indutancia      = 2.5000e-03 H
	Capacitancia    = 6.2500e-07 F
	Duty-cycle      = 0.5000
	Corrente media  = 2.0000 A
	Corrente maxima = 2.5000 A
	Corrente minima = 1.5000 A
\end{lstlisting}

\begin{lstlisting}[caption={Exemplo 2 - DCM.},label={lst:item5b}]
	>> projeto_boost(100,200,10e3,200,110,1)
	
	-------------------------------
	Projeto do Conversor Boost
	-------------------------------
	D = 0.5000
	K = 0.1136
	Kcrit = 0.1250
	
	Modo de operacao: DCM
	
	Indutancia = 1.2500e-03 H
	Duty-cycle = 0.5000
\end{lstlisting}
